\documentclass[10pt,a4paper]{amsart}
\usepackage[margin=19mm]{geometry}
\usepackage{amsmath,amssymb,mathtools}
\usepackage[scr=rsfs,cal=euler]{mathalfa}
\usepackage{xfrac}
\usepackage[unicode,hidelinks,bookmarks=false]{hyperref}
\newcommand{\ie}{i.e.\ }
\newcommand{\cf}{cf.\ }
\newcommand{\resp}{respectively\ }
\newcommand{\Subsection}{\subsection}
\providecommand{\nobreakdash}{\nobreak\hbox{-}}
\setlength{\emergencystretch}{2em}
\multlinegap0pt
\usepackage[shortlabels]{enumitem}
\usepackage{amssymb}
\usepackage{mathtools}
\usepackage{bm}
\newtheorem{theorem}{Theorem}
\usepackage{cleveref}
\crefname{theorem}{Theorem}{Theorems}
\newcommand{\AI}{\mathrm{AI}}
\newcommand{\R}{\mathbb R}
\newcommand{\SPD}{\mathbb S_{++}}
\newcommand{\argminop}{\operatorname*{arg\,min}}
\newcommand{\dai}{d_{\AI}}
\newcommand{\fro}[1]{\left\lVert #1\right\rVert_F}
\newcommand{\norm}[1]{\left\lVert #1\right\rVert}
\title{Information-Induced Training Geometry:\\ Exact Reduction, Canonical Completion, and Structured Expressivity}
\author{Zavier Li}
\date{Source: arXiv:2609.12991v1 (CC BY 4.0)}
\begin{document}
\maketitle

\noindent\emph{Source and licence.} Information-Induced Training Geometry:\\ Exact Reduction, Canonical Completion, and Structured Expressivity. Version arXiv:2609.12991v1 (CC BY 4.0), distributed by arXiv under the Creative Commons Attribution 4.0 International licence, http://creativecommons.org/licenses/by/4.0/, as recorded in the arXiv licence metadata for this exact version and shown on the abstract page https://arxiv.org/abs/2609.12991v1. Full licence notice: \texttt{licenses/arxiv-2609.12991-CC-BY-4.0.txt}.\par\smallskip

\noindent\textit{Editorial scope.} Counted unit: the article's theorem thm:spd-compression (source0.tex lines 608--673). The article's complete original proof of it is delivered with it (source0.tex lines 1685--1798, one \texttt{proof} environment of 114 lines, which the article places away from the statement). No mathematics is written, repaired or re-derived: the statement and the proof are the article's own lines, in the article's own notation, and no retained line is substituted or re-flowed.  Retained uncounted as the source-local context the proof needs: lines 587--595.  Nothing was omitted from any retained span, so the package carries no omission record.  No retained line contains a citation, so the wrapper carries no bibliography and no bibliography entry.  No retained line contains a figure environment, an image-inclusion command or drawing code, so no image is delivered and the package declares no asset dependency.  Editorial formatting only, in addition: an AMS \texttt{amsart} preamble that restates the article's own declarations and macros used by the retained lines, namely the environments \texttt{theorem} on separate counters, together with the standard packages those lines need.  The retained environments print in this document as Theorem 1 in order of appearance, whereas the article prints the same items with the numbering of its own full document.  The complete native source of the article as distributed in the arXiv e-print for this exact version is bundled unchanged as \texttt{sources/210167/source0.tex}.
\begin{align}
  \Psi_{P_0,A}(R)
  &=P_0^{1/2}
  \left[I+U_0(R_0^{-1/2}RR_0^{-1/2}-I)U_0^\top\right]
  P_0^{1/2}
  \nonumber\\
  &=P_0+P_0AR_0^{-1}(R-R_0)R_0^{-1}A^\top P_0.
  \label{eq:canonical-completion}
\end{align}

\begin{theorem}[Complete SPD information-compression theorem]
\label{thm:spd-compression}
The map \(\pi_A\) and the section \(\Psi_{P_0,A}\) satisfy the following.
\begin{enumerate}[label=(\Alph*),leftmargin=*]
  \item \emph{Contraction and submetry.}  For all \(P,Q\in\SPD^d\),
  \begin{equation}
    \dai(A^\top PA,A^\top QA)\le\dai(P,Q),
    \label{eq:spd-contraction}
  \end{equation}
  and for every \(r\ge0\),
  \begin{equation}
    \pi_A(\overline B_{\AI}(P_0,r))
    =\overline B_{\AI}(R_0,r).
    \label{eq:spd-ball-submetry}
  \end{equation}
  \item \emph{Unique canonical completion.}  For every \(R\in\SPD^m\),
  \begin{equation}
    \boxed{
    \Psi_{P_0,A}(R)
    =\argminop_{A^\top PA=R}\dai(P_0,P),
    \qquad
    \dai(P_0,\Psi_{P_0,A}(R))=\dai(R_0,R).}
    \label{eq:spd-shortest-completion}
  \end{equation}
  The minimizer is unique.
  \item \emph{Horizontal geometry.}  The section \(\Psi_{P_0,A}\) is an
  isometric totally geodesic embedding.  The map \(\pi_A\) is a Riemannian
  submersion; at \(P\), the inverse of
  \(D\pi_A(P)|_{\mathcal H_P}\) is
  \begin{equation}
    h\longmapsto D\Psi_{P,A}(A^\top PA)[h].
    \label{eq:spd-horizontal-lift}
  \end{equation}
  Hence \(\pi_A\) is split-Hadamard and visible AIRM geodesics have unique
  horizontal ambient lifts.
  \item \emph{Invisible-fiber minimax geometry.}  If \(m<d\), the fiber
  \(\mathcal F_A(R)=\{P:A^\top PA=R\}\) has infinite AIRM diameter.  With
  \(P_\star=\Psi_{P_0,A}(R)\) and
  \[
    \mathcal F_A(R;B)
    =\mathcal F_A(R)\cap\overline B_{\AI}(P_\star,B),
  \]
  one has
  \begin{equation}
    \inf_{Q\in\SPD^d}\sup_{P\in\mathcal F_A(R;B)}\dai(Q,P)=B,
    \label{eq:spd-minimax}
  \end{equation}
  and \(P_\star\) is the unique minimax center.
  \item \emph{Exact radial decision reduction.}  Let \(\mathcal A\) be any
  family of admissible \((P_0,A)\), allowing different visible dimensions,
  let \(\rho\) be nondecreasing, and let \(\varphi\) be any extended-real
  visible functional.  Then
  \begin{align}
    &\inf_{\substack{(P_0,A)\in\mathcal A\\P\in\SPD^d}}
    \left\{\rho(\dai(P_0,P))+\varphi(P_0,A,A^\top PA)\right\}
    \nonumber\\
    &\quad=
    \inf_{\substack{(P_0,A)\in\mathcal A\\R\in\SPD^m}}
    \left\{\rho(\dai(A^\top P_0A,R))+\varphi(P_0,A,R)\right\}.
    \label{eq:spd-decision-reduction}
  \end{align}
  Every reduced minimizer has the canonical full lift
  \(P=\Psi_{P_0,A}(R)\).  If \(\rho\) is strictly increasing on its finite
  domain, every finite-valued full minimizer is canonical.
\end{enumerate}
\end{theorem}

\begin{proof}[Proof of \cref{thm:spd-compression}]
Fix \(P\in\SPD^d\), write \(R=A^\top PA\), and let \(H=H^\top\) be a
tangent at \(P\).  Define
\[
  C=P^{1/2}AR^{-1/2},
  \qquad
  X=P^{-1/2}HP^{-1/2}.
\]
Then \(C^\top C=I_m\) and
\[
  R^{-1/2}A^\top HA R^{-1/2}=C^\top XC.
\]
Complete \(C\) to an orthogonal matrix \([C,N]\).  The exact block identity
is
\begin{equation}
  \fro X^2
  =\fro{C^\top XC}^2
  +2\fro{N^\top XC}^2
  +\fro{N^\top XN}^2.
  \label{eq:block-frobenius-identity}
\end{equation}
Therefore
\[
  \norm{D\pi_A(P)[H]}_{R,\AI}
  =\fro{C^\top XC}
  \le\fro X
  =\norm H_{P,\AI}.
\]
Integrating along curves proves global contraction
\eqref{eq:spd-contraction}.

Congruence by \(P_0^{-1/2}\) in the ambient cone and by \(R_0^{-1/2}\) in
the visible cone reduces the completion problem to
\[
  \min\{\dai(I,X):U_0^\top XU_0=\widetilde R\},
  \qquad
  \widetilde R=R_0^{-1/2}RR_0^{-1/2}.
\]
The candidate
\[
  X_\star=I+U_0(\widetilde R-I)U_0^\top
\]
is positive definite, has compression \(\widetilde R\), and satisfies
\[
  \dai(I,X_\star)=\fro{\log\widetilde R}=\dai(I,\widetilde R).
\]
Contraction proves it is a minimizer.

For uniqueness, suppose feasible \(X\) attains equality and write
\(L=\log X\).  The compressed curve
\(Y_t=U_0^\top\exp(tL)U_0\) joins \(I\) to \(\widetilde R\).  Hence
\[
  \dai(I,\widetilde R)
  \le\operatorname{Len}(Y)
  \le\fro L
  =\dai(I,X)
  =\dai(I,\widetilde R).
\]
Equality of the continuous speed inequality at \(t=0\), together with
\eqref{eq:block-frobenius-identity}, forces
\(N^\top LU_0=0\) and \(N^\top LN=0\).  Thus
\(L=U_0L_0U_0^\top\), and the endpoint constraint gives
\(\exp(L_0)=\widetilde R\).  Therefore \(X=X_\star\).  Undoing congruence
gives \eqref{eq:canonical-completion} and
\eqref{eq:spd-shortest-completion}.

For \(R,S\in\SPD^m\), the normalized lifts are
\(\widetilde R\oplus I\) and \(\widetilde S\oplus I\).  Their ambient AIRM
distance is the visible AIRM distance, and their geodesic keeps the identity
block fixed.  Thus the section is isometric and totally geodesic.  The ball
identity follows by combining contraction with the equal-distance lift.

At a general \(P\), let
\(L_P(h)=D\Psi_{P,A}(A^\top PA)[h]\).  The section identity and isometry give
\[
  D\pi_A(P)[L_P(h)]=h,
  \qquad
  \norm{L_P(h)}_{P,\AI}=\norm h_{A^\top PA,\AI}.
\]
If \(V\in\ker D\pi_A(P)\), then \(L_P(h)+sV\) is another lift of \(h\).
Infinitesimal contraction gives
\(\norm h^2\le\norm{L_P(h)+sV}^2\) for all \(s\), with equality at zero.
Differentiation shows \(L_P(h)\perp V\).  Surjectivity and dimension counting
identify the image of \(L_P\) with the full horizontal space.  This proves
the Riemannian-submersion and horizontal-lift statements.  The images of the
sections through all \(P\) are complete totally geodesic maximal horizontal
integral leaves; hence the horizontal distribution is integrable and
\(\pi_A\) is split-Hadamard.

For the minimax statement, work in the normalized splitting and choose
symmetric \(H\in\R^{(d-m)\times(d-m)}\) with \(\fro H=1\).  The curve
\[
  X_s=\widetilde R\oplus\exp(sH)
\]
stays in the visible fiber and satisfies
\(\dai(X_0,X_s)=|s|\).  This proves infinite diameter.  The two points
\(X_{-B},X_B\) lie in the bounded fiber and are distance \(2B\) apart, so
every center has worst-case distance at least \(B\).  The canonical point
\(X_0\) has worst-case distance at most \(B\).  If another point attained
\(B\), equality in the triangle inequality between the endpoints would make
it their midpoint.  AIRM is uniquely geodesic, so the midpoint is uniquely
\(X_0\).

Finally, for any admissible \((P_0,A)\), every \(P\) with
\(R=A^\top PA\) obeys
\[
  \dai(P_0,P)\ge\dai(A^\top P_0A,R),
\]
and the canonical lift realizes equality.  Monotonicity of \(\rho\) proves
both inequalities in \eqref{eq:spd-decision-reduction}.  The same construction
lifts every reduced minimizer.  Under strict monotonicity, replacing a
noncanonical full point by the strictly nearer canonical point lowers the
objective, proving the final correspondence.
\end{proof}

\end{document}
